\begin{claim}{B.38}{Proved}{The free positive involution and the six negative involutions in genus three}
There are exactly twelve conjugacy classes of mixed finite subgroups
of $\mathcal T_{3,0}$ whose positive subgroup is a free involution.
Each is a four-group.  Their quotient signatures and numbers are
\[
\begin{array}{c|ccccc}
\text{underlying quotient}&S_{0,3}&S_{1,1}&N_{1,2}&N_{2,1}&N_{3,0}\\\hline
\text{number of classes}&4&1&4&2&1 ,
\end{array}                                             \tag{B.38.1}
\]
where $S_{h,b}$ is orientable of genus $h$ with $b$ mirror boundary
components, and $N_{q,b}$ is nonorientable with $q$ crosscaps and
$b$ mirror boundary components.  There are no cones or corners.
There are also exactly six mixed classes whose positive subgroup
is trivial: the negative involutions with quotient signatures
\[
 S_{0,4},\quad S_{1,2},\quad
 N_{1,3},\quad N_{2,2},\quad N_{3,1},\quad N_{4,0}.
                                                        \tag{B.38.2}
\]
All eighteen classes, their exact monodromies, complete section
posets, $\Phi$, stabilizers, centralizers and multipliers appear in
\path{type_3_0_free_positive_and_anti_catalogue.csv}
and its \texttt{sections.json} companion.
\end{claim}
\begin{proof}
Let $X$ be the closed orientable genus-three surface.  First let
$G^+=\langle a\rangle$ have order two and act freely, and let
$G$ contain negative elements.  Then $|G|=4$.
A group of order four is cyclic or a four-group: an element of
order four gives the first case; otherwise $(xy)^2=1$ for all
$x,y$, which gives commutativity and the second case.
There is no nontrivial positive point stabilizer.
A negative point stabilizer is generated by a reflection, so
$Q=X/G$ has no cone or corner points.  Its orbifold Euler
characteristic is $-4/4=-1$.  The complete list of compact
underlying surfaces with this Euler characteristic is exactly
(B.38.1).

The cyclic possibility cannot occur.  A negative element of
order four has no fixed point, since a negative local finite
stabilizer element is a reflection of order two.  Its square
is free by assumption, so the whole action would be free.
The quotient would then be $N_{3,0}$.
Its fundamental group has the presentation
\[
 \langle u_1,u_2,u_3\mid u_1^2u_2^2u_3^2=1\rangle .
\]
An orientation-compatible monodromy to $C_4=\Z/4$ sends each
crosscap generator to an odd residue.  Each square has image
$2$, so the relator has image $6=2\pmod4$, a contradiction.
The program checks all eight assignments of odd residues.
Thus $G$ is the four-group
$V_4=\mathbf F_2^2$, with orientation given by its second
coordinate, and $a=(1,0)$.

We classify all four-group monodromies, including their full
topological equivalence relation.  The mirror orbifold group
of an orientable $S_{h,b}$ has generators
$a_i,b_i,e_j,c_j$ and relations
\[
 c_j^2=1,\qquad[e_j,c_j]=1,\qquad
 \prod_i[a_i,b_i]\prod_je_j=1.                         \tag{B.38.3}
\]
The negative $c_j$ is the reflection on the $j$th boundary,
and the positive $e_j$ traverses that boundary.
For $N_{q,b}$ replace the handles by crosscap generators
$u_i$, and replace the last relation by
\[
                         \prod_i u_i^2\prod_j e_j=1.  \tag{B.38.4}
\]
The orientation character $\varepsilon$ is zero on handles
and $e_j$, and one on crosscaps and $c_j$.
An orientation-compatible epimorphism to $V_4$ is exactly
\[
 q_\alpha:\Pi_Q\longrightarrow V_4,\qquad
 q_\alpha(x)=(\alpha(x),\varepsilon(x)),
 \quad
 \alpha\in H^1(\Pi_Q;\mathbf F_2)
       \setminus\{0,\varepsilon\}.                   \tag{B.38.5}
\]
In coordinates write the handle or crosscap values as $\eta_i$,
the boundary values as $\beta_j=\alpha(e_j)$, and the reflection
values as $\gamma_j=\alpha(c_j)$.
The only linear relation is
$\sum_j\beta_j=0$.
The two automorphisms of $V_4$ preserving orientation identify
$\alpha$ with $\alpha+\varepsilon$.

Every map (B.38.5) is injective on local reflection groups.
Its kernel is therefore torsion-free; one proves this either
from the universal hyperbolic orbifold cover or by cutting along
arcs to its reflection polygon, where every finite stabilizer
is a displayed local reflection group.
The kernel lies in the orientation kernel, and its Euler
characteristic is $4(-1)=-4$.  It is consequently the fundamental
group of a closed orientable genus-three surface.
Conversely every action under consideration yields exactly
this construction.  Two such actions are conjugate precisely
when an orbifold homeomorphism of $Q$ carries one monodromy
kernel to the other, or equivalently when it carries the
monodromy to the other after an automorphism of $V_4$.
This follows by lifting paths in the two covers, with the
local reflection charts giving the extension across mirrors.
Thus the following finite cohomology-orbit calculation is an
exact classification of actions.

We prove the entire mapping class action needed in this
calculation.  A boundary permutation permutes the pairs
$(\beta_j,\gamma_j)$.  The interior homology action preserves
the mod-two intersection pairing, its radical generated by
boundary classes, and the orientation character.
It introduces no reflection coordinate: an interior loop
is mapped to an interior loop, whereas each reflection generator
is mapped to a conjugate of the reflection at the corresponding
boundary.  Thus the $\gamma_j$ are affected only by boundary
permutations and by the global replacement
$\alpha\mapsto\alpha+\varepsilon$.

For $S_{0,3}$ the interior homology is generated by its boundary
classes, with their single sum relation, so boundary permutations
are the whole action.  When $\beta=(0,0,0)$, surjectivity leaves
one nonconstant reflection-color pattern up to permutation and
global color reversal.  When $\beta$ has two ones, distinguish
the boundary with value zero and use global color reversal to
make its reflection value zero.  The two remaining reflection
values have weights zero, one or two.  This gives $1+3=4$ classes.

For $S_{1,1}$ the boundary value is zero.
Global color reversal sets its reflection value to zero.
The remaining nonzero pair of handle values is a single orbit
under $\mathrm{Sp}_2(\mathbf F_2)$, generated by the two Dehn
twists in a pair of basis curves.  This gives one class.

For $N_{1,2}$ write $u$ for its one-sided core and $e$ for
either mod-two boundary class.  Its intersection form has
$u\cdot u=1$, with $e$ in the radical.
An automorphism preserving the boundary classes and the
orientation character sends $u$ either to $u$ or to $u+e$.
Both possibilities are realized.  Push one boundary disc once
around the one-sided core and return it as an unparametrized
disc; its boundary orientation may reverse.  A core arc crossing
the track once acquires one boundary circuit.  Thus, modulo two,
this homeomorphism sends $u$ to $u+e$ and leaves both reflection
coordinates unchanged.  The boundary exchange is realized away
from the core.  These maps generate every permitted action on
the displayed homology and boundary data.
After global color reversal sets $\eta=\alpha(u)$ to zero,
the case $\beta=(0,0)$ has two nonzero reflection patterns up
to exchange.  The case $\beta=(1,1)$ has two patterns: the
boundary push followed by global reversal interchanges
$\gamma=(0,0)$ and $(1,1)$, while $(0,1)$ and $(1,0)$ are
exchanged by the boundary permutation.  This gives four classes.
The boundary push is necessary for this count.

For $N_{2,1}$ the boundary value is zero.
Set its reflection value to zero by global reversal.
The intersection form on the two crosscap classes is the
identity matrix over $\mathbf F_2$; its orthogonal group has
exactly the identity and the exchange of those two classes.
They are realized by exchanging the two crosscaps.
The nonzero crosscap vectors therefore have two orbits,
represented by $(1,0)$ and $(1,1)$.

For $N_{3,0}$ the mod-two intersection form is the identity
matrix in three crosscap classes, and its characteristic
vector is $(1,1,1)$.  An orthonormal basis cannot contain that
characteristic vector, since it pairs to one with every vector
of odd self-intersection.  Its members must therefore be the
three standard basis vectors.  The full orthogonal group is
their permutation group, realized by crosscap permutations.
Global addition of $\varepsilon=(1,1,1)$ exchanges weights
one and two, leaving one orbit after $0,\varepsilon$ are
excluded.  This proves all five numbers in (B.38.1).

We now compute the sections explicitly.  A mirror boundary
with data $(\beta_j,\gamma_j)$ has local reflection
$c_j=(\gamma_j,1)$ and boundary monodromy
$e_j=(\beta_j,0)$.
Its preimage is described by the two cosets
$V_4/\langle c_j\rangle$; one trip around the boundary
permutes these cosets by $e_j$.
Consequently it gives two fixed ovals when $\beta_j=0$,
and one when $\beta_j=1$.
The two-oval case has an oval stabilizer of order two,
$\langle c_j\rangle$, and the one-oval case has stabilizer
$V_4=\langle c_j,e_j\rangle$.
Both are one $V_4$-orbit of section classes.
There are no point sections because the positive involution
is free, and there are no full sections because $V_4$ cannot
fix a point without giving a positive local stabilizer.
Every reflection section is maximal, has centralizer $\Z$
in the surface group, and has multiplier $1/2$ on its
nonidentity element.  Thus
\[
 |\Phi|=\#\{\text{nontrivial partial sections}\}
       =2b-\#\{j:\beta_j=1\},
 \qquad |\Phi/V_4|=b.                                 \tag{B.38.6}
\]
The program constructs the actual finite coset graph, with
three subdivisions per boundary traversal, computes its
components and the group action, and verifies every count
and stabilizer in this formula.

Finally a mixed group with trivial positive subgroup is $C_2$.
Its quotient has orbifold Euler characteristic $-2$ and no
cones or corners.  An orientable quotient must have mirror
boundary; otherwise its orientation character would be
trivial and could not give the required negative covering.
Solving $2-2h-b=-2$ with $b>0$, or $2-q-b=-2$ with $q>0$,
gives precisely (B.38.2).
For each quotient the monodromy is forced: it is the orientation
character, taking every crosscap and local reflection to the
nonidentity element and every handle and boundary traversal
to zero.  Its orientation double is connected, torsion-free,
and has Euler characteristic $-4$.  The classification of
compact quotient surfaces and path lifting prove that each
signature gives one conjugacy class and that distinct signatures
give distinct classes.
There are $b$ fixed ovals, hence $b$ full sections, $b$
nontrivial partial sections, and $b$ maximal classes, each
fixed by the whole $C_2$ and with surface-group centralizer
$\Z$.  The case $N_{4,0}$ has none.  These are the six
genus-three real types of B.16 in an executed form.

The code \texttt{catalogue\_genus3\_free\_positive.py} enumerates
every cohomology vector in (B.38.5), applies exactly the
proved generators of the equivalence relation, checks
surjectivity, all quotient relations, orientation compatibility,
and the covering Euler characteristic, and records every
orbit.  Its second calculation constructs all six orientation
doubles in (B.38.2).  This supplies the complete eighteen-row
catalogue and its section data.
\end{proof}
